\scanpage{25}
isolé $t\in\Omega$ tel que $U_{\varphi_{\{t\}}}E=EU_{\{t\}}=E$, \deleted{\unsure{par}} \inserted{d'où}

$A$ est le \deleted{\unsure{somme} directe} \inserted{produit} de deux sous-algèbres $A_0$ et $A_{2}$, où $A_1$ est un \uline{facteur}, ayant de un projecteur minimal $E$. \deleted{Cela prouve que} \inserted{dans $A$,}
\deleted{\inserted{l'algèbre de ses \unsure{éléments} propres} $A$ est le ``produit'' (au sens des produit\unsure{s}) de \unsure{deux} algèbres $A_0$ et $A_1$, où le $A_{\unsure{1}}$ est}
il est de type I, \deleted{$T$ induit} donc isomorphe à un $L(\mathcal H)$, $T$ induit sur $A_1$ \inserted{un \deleted{facteur} \inserted{\unsure{nombre}} positif} la trace usuelle, et $T(E)<+\infty$, \deleted{Enfin, si} d'où $E\in\underline a\subset\underline b$.

\deleted{Enfin, si \inserted{alors} $E$ est cet}
\inserted{Soit alors $\underline c=C_0(M)$, algèbre} \inserted{auto-adjointe} commutative maximale \inserted{contenue dans $\underline b$} contenant $E$, \deleted{alors} \deleted{on sait que} $E$ s'identifie à la fonction caractéristique d'un ensemble à la fois ouvert et fermé de $M$, qui se réduit nécessairement à un pt ($E$ étant minimal), et $\unsure{\mu|\nu}(\{t\})>0$. \deleted{\illegible{2}}

Une telle algèbre est dite \deleted{de type I pur}.
\deleted{Ainsi on a} \deleted{De plus} \deleted{Ainsi, pour que la \unsure{mesure soit}}
\deleted{
Ainsi en définitive \inserted{comme $\underline c=C_0(M)$, \deleted{$\underline b$} \unsure{de}} \deleted{\unsure{compte} dans un} la mesure $\unsure{\mu|\rho}$ sur $M$ soit non atomique, il faut et il suffit que $A$ n'ait pas de projecteur minimal. Dans le cas général,

Soit $I$ l'ens des $i\in\unsure{\mathcal C|\mathcal E}$ \inserted{isolés} du spectre $\Omega$ du centre $C$ de $A$, tels que les sous-algèbres $A_i=U_{\{i\}}AU_{\{i\}}$ de $A$ \deleted{qui sont des} (facteurs) sont de \uline{type I}, i.e. admettent un projecteur minimal. On sait alors
}
\mspage{26}
